\section{Introduction}
\label{sec:intro}

Reconnaissance precedes damage in attacks on power-grid control systems. Before an adversary can
open a breaker or forge a measurement that passes the state estimator's bad-data
detection~\cite{liufdia2009}, it first learns which devices are on the network and how they behave. In December 2015, remote intruders mapped a Ukrainian distribution utility, then used that
knowledge to open breakers and cut power to roughly 225{,}000 residents~\cite{ukraine2015}. The
information a passive observer can extract from unencrypted control traffic is therefore an early and
consequential step in the attack chain, not a theoretical concern.

A passive observer does not need to decode a DNP3 payload to fingerprint an outstation. Formby
et~al.~\cite{formby2016} show that cyber-physical devices are identifiable from cross-layer response
timing and physical-operation timing, features that survive the absence of payload decoding. Two
wire-visible features carry this signature for DNP3. The cross-layer response time (CLRT) is the
interval between the transport-layer ACK and the application-layer response; it reflects the
outstation's internal processing and is stable across sessions. The segment shape is the vector of
TCP segment sizes into which the outstation divides a response; because a DNP3 outstation segments
deterministically, the shape repeats identically. On our SEL-751 relay, a native read answers with a
single 49-byte segment and a CLRT median of 1.3~ms, while a control select answers at 2.1~ms, so the
two transaction classes are separable by timing alone. Prior traffic-analysis work shows that such
size-and-timing side channels recover content even on encrypted sessions~\cite{sirinam2018}.

Normalizing these features in the network is desirable for two reasons. \emph{First}, an in-network
defense needs no change to the master or the vendor-supplied outstation, which a legacy substation
cannot modify. \emph{Second}, it suppresses the reconnaissance signal before an attack executes,
rather than detecting the attack after the fact. The deployment is, however, constrained on both
sides. \emph{First},
transport encryption is usually absent on legacy DNP3 and cannot be assumed, and encryption alone
would not hide the timing and segment-count features that survive tunneling~\cite{sirinam2018}.
\emph{Second}, the outstation firmware is fixed and vendor-supplied, so a defense cannot change how
the device builds its frames, and any modification of the response bytes breaks the per-block CRCs the
master validates~\cite{dnp3std}. \emph{Third}, generic traffic-analysis defenses do not fit: per-packet
padding leaves exploitable size and count features~\cite{dyer2012}, morphing one class into another
assumes control of the sending endpoint~\cite{wright2009}, and website-fingerprinting defenses such as
adaptive padding target browser traffic over an anonymity network with host cooperation~\cite{juarez2016}.
\emph{Last}, host-based delay or a controller-mediated release changes the endpoints or adds a
control-plane fast path that a legacy substation deployment cannot accommodate.

There exists a gap for a defense that normalizes both the response timing and the eligible segment
shape of DNP3 traffic in the network while preserving protocol validity and requiring
no endpoint change. We address this gap with an in-network obfuscation layer that runs entirely in the
data plane of one Intel Tofino-1 switch placed at the outstation edge. The defense holds each response
in strict-priority switch queues and releases its ACK and response at fixed policy offsets from the
request arrival, so the master-visible CLRT becomes a constant. It then replicates the eligible
49-byte response and carves it at an existing DNP3 CRC-block boundary into a fixed $[28,21]$ segment
vector, so no DNP3 byte is modified and the master reassembles the identical response. We refer to this
timing-and-segmentation path as response shaping.

The defense exposes selectable modes at runtime. A passthrough mode forwards traffic natively. Timing
modes hold the ACK, the response, or both to their policy offsets; the combined mode, which holds both,
is the one that normalizes the CLRT and drives the response-shaping path above. A separate mode
protects control commands.

The control-command mode is necessary because the Select-before-Operate exchange exposes an
operation-timing fingerprint: the delay from an OPERATE to its confirmation reflects the relay's
actuation. Response shaping alone cannot hide this delay, because releasing the OPERATE on a fixed
schedule would still reveal it. The control-command mode blocks the OPERATE and is designed to release
it once at a secret internal delay while anchoring the master-visible ACK and echo to the original
arrival time, so the observable gap does not vary with the hidden delay. Response shaping and the
control-command mode cannot share one strict-priority queue ladder, so the defense runs them on two
independent switch scheduling domains, both internal to the same Tofino and the same program;
Section~\ref{sec:impl} explains the scheduling conflict that forces the split.

We evaluate the defense on a testbed with a physical master, one Tofino-1, and a physical SEL-751 relay,
comparing native and defended traffic captured on the same binary under a runtime mode change. On the
master-facing link, the defended CLRT is a fixed 4.001~ms (standard deviation 0.02~ms) for both read
and select, in place of the native 1.3~ms and 2.1~ms; every one of 1{,}280 defended eligible responses
presents the $[28,21]$ segment vector with a sequence-contiguous, CRC-valid, checksum-valid 49-byte
reconstruction and no unsplit escape. The control-command gap stays at 4.00~ms across the evaluated release delays, and a classifier that separated read from select by CLRT with 0.59 balanced accuracy
on native traffic drops to the 0.50 chance baseline on defended traffic. The guarded control test
permitted only isolated output points and refused a breaker-close index, and all 32 relay outputs
remained open throughout.

The evaluation is bounded. We test one outstation, so the defense replaces one device's CLRT and
segment-shape signature with a policy rather than establishing indistinguishability among devices. The
relay-facing link was not captured, so the single release at the secret delay is verified by the
implementation model but not observed on the wire, and we do not claim exactly-once delivery to the
relay. The defense normalizes the observed segment vector, not the total 49-byte payload length.

To the best of our knowledge, this is the first in-network defense to normalize both the response
timing and the segment shape of a live industrial-control protocol while preserving its bytes and
CRCs. This paper makes the following contributions.
\begin{itemize}
\item \textbf{A hardware timing primitive for absolute-offset response shaping.} The response-shaping
path holds the ACK and response in strict-priority reservoirs and releases them at fixed offsets from
the request timestamp, normalizing the master-visible CLRT to a constant, with selectable timing modes.
\item \textbf{Deterministic in-network segment-shape normalization.} We use packet replication and
egress carving to emit an eligible response as a fixed $[28,21]$ segment vector on an existing DNP3
CRC-block boundary, preserving every DNP3 byte and CRC.
\item \textbf{A control-command timing mode.} The control-command mode blocks the OPERATE and anchors
the master-visible ACK and echo to the request arrival, and we identify why the two mechanisms require
two independent switch scheduling domains.
\item \textbf{A one-program implementation and a bounded hardware evaluation.} We implement the defense
as one P4 program that fits a single Tofino-1 ingress pipe within twelve stages, and we evaluate it on
a physical SEL-751 with a reproducible capture-to-claim analysis, an explicit safety result, and a
stated evidence boundary.
\end{itemize}
